The command \textit{HeteroAgentStationaryEqm\_{}Case1} (and \textit{\_{}Case2}) allows for easy calculation of the price (vector) associated with the general equilibrium of
standard Bewley-Huggett-Aiyagari models\footnote{General equilibrium incomplete market heterogeneous agent models with idiosyncratic, but no aggregate, shocks.}$^{,}$\footnote{In
principle one could use these codes to solve fixed-point problems around value function problems more generally, but this would be more advanced.}.

As well as the standard inputs required for the Value function iteration command the other inputs needed are related to evaluating the general equilibrium
conditions.

By default the command will use matlabs \href{http://www.mathworks.com/help/matlab/ref/fminsearch.html}{fminsearch} command to solve the fixed-point problem represented by the general equilibrium requirement.
Alternatively the user can set $n\_{}p$ to be non-zero in which case the general equilibrium condition will be evaluated on a grid for prices. This can
be useful if you which to check for the possible existence of multiple equilibria and has the advantage of possessing well understood convergence
properties \citep*{KirkbyHeteroErrors}.

In the notation of the toolkit this is any problem for which the competitive equilibrium can be written as
\begin{definition} \label{defn:HeteroAgentCE}
A \textit{Competitive Equilibrium} is an agents value function $V_{p}$; agents policy function $g$;
vector of prices $p$; measure of agents $\mu_{p}$;
such that
\begin{enumerate}
 \item Given prices $p$, the agents value function $V_p$ and policy function $g_p$ solve the agents problem
   given by a Case 1 value function (as described in Section \ref{sec:VFI_Case1}).
 \item Aggregates are determined by individual actions: $B_p=\mathcal{A}(\mu_p)$.
 \item Markets clear (in terms of prices): $\lambda(B_p,p)=0$.
 \item The measure of agents is invariant:
      \begin{equation}
          \mu_p(a,z)=\int \int \left[ \int 1_{a=g^{a'}_p(\hat{a},z)}\mu_p(\hat{a},z) Q(z,dz')\right] d\hat{a}dz
      \end{equation}
\end{enumerate}
\end{definition}
the $p$ subscripts denote that these objects depend on the price(s) $p$.

\subsection{Preparing the model}
To use the toolkit to solve problems of this sort the following steps must first be made.
\begin{enumerate}
 \item Create all of the variables that are needed for the value function problem. These are exactly the same as those used for
   $ValueFnIter\_{}Case1$ and the exact same notation applies (See Section \ref{sec:VFI_Case1}). \\
    (Namely \textcolor{RoyalBlue}{V0, n\_{}d, n\_{}a, n\_{}z, d\_{}grid, a\_{}grid, z\_{}grid, pi\_{}z,} \\
    \indent \textcolor{RoyalBlue}{ReturnFn, Parameters, DiscountFactorParamNames, ReturnFnParamNames, [vfoptions]})
 \item Define functions for any aggregate variables that are needed by the general equilibrium equation in terms of the micro-level variables, these
   are evaluated as the integral of the function over the stationary agents distribution. Store all these functions together. \\
   eg., if the only aggregate variables is the integral of the endogenous state, then \\
   \textcolor{RoyalBlue}{FnsToEvaluateFn\_{}1 = @(aprime\_{}val,a\_{}val,s\_{}val;} \\
   (You always need the $@(...)$ defined either in this way, or with $@(d\_{}val,...)$ where appropriate. You also need to include parameters where appropriate, done in same way
   as parameters are passed to the $GeneralEqmEqn$ below.) \\
   Define the names of any parameters used to evaluate the aggregate variables. \\
   \textcolor{RoyalBlue}{FnsToEvaluateParamNames(1).Names=\{\};} \\
   (In our example there were none, but an example might be that you want to evaluate tax revenue, which is a tax rate parameter times an endogenous state variable, then the tax rate
    parameter name will be given here.) \\
   You then store all of these together: \\
   \textcolor{RoyalBlue}{FnsToEvaluate=\{FnsToEvaluateFn\_{}1\};}
   (For models with multiple FnsToEvaluateFns simply create more in same manner and store them all in $FnsToEvaluate$.)
 \item Define functions for how to calculate general equilibrium conditions for which we need to find general equilibrium. Then store these
    functions together. \\
    Create General Equilibrium equation: \\
    \textcolor{RoyalBlue}{GeneralEqmEqn\_{}1 = @(AggVars,p,param1, param2) p-param1*AggVars(1)\^{}param2} \\
    Give the names of the parameters: \\
    \textcolor{RoyalBlue}{GeneralEqmEqnParamNames(1).Names = $\{$'param1', 'param2'$\}$} \\
    ($param1$ must exist in $Parameter.param1$; similarly for $param2$. '$param1$' is the name of the parameter, e.g. '$alpha$') \\
    You then store all of these together: \\
    \textcolor{RoyalBlue}{GeneralEqmEqns={GeneralEqmEqn\_{}1};} \\
    (For models with multiple General Equilibrium Equations simply create more in same manner and store them all in $GeneralEqmEqns$.)
 \item Define the names of the prices which are being used to calculate general equilibrium equations. General equilibrium will be where the inputed
   values of these correspond to those that lead the general equilibrium equations to be zero-valued. \\
   \textcolor{RoyalBlue}{GEPriceParamNames=\{'r'\};} \\
   (Technically, in terms of finding the general equilibrium these could be any quantity, not just prices.)
\end{enumerate}
That covers all of the objects that must be created, the only thing left to do is simply call the heterogeneous agent general equilbrium code. \\
\textcolor{RoyalBlue}{[p\_{}eqm,p\_{}eqm\_{}index,GeneralEqmConditions] =} \\
\indent \textcolor{RoyalBlue}{  HeteroAgentStationaryEqm\_{}Case1(V0, n\_{}d, n\_{}a, n\_{}z, n\_{}p, pi\_{}z,d\_{}grid, a\_{}grid, s\_{}grid,} \\
\indent \indent \textcolor{RoyalBlue}{   ReturnFn, FnsToEvaluate, GeneralEqmEqns, Parameters, DiscountFactorParamNames, } \\
\indent \indent \textcolor{RoyalBlue}{   ReturnFnParamNames, FnsToEvaluateParamNames, GeneralEqmEqnParamNames, } \\
\indent \indent \textcolor{RoyalBlue}{   GEPriceParamNames, [heteroagentoptions], [simoptions], [vfoptions]);}

Once run $p\_{}eqm$ will contain the general equilibrium price vector. $GeneralEqmConditions$ will be a vector containing the values for each of the market
clearance conditions (general equilibrium is where these equals zero).
$p\_{}eqm\_{}index$ will simply take the value $nan$, unless you are using the
(non-default) grid on prices options, see 'Some further remarks' below.

\subsection{Some further remarks}
\begin{itemize}
 \item The $Case2$ command is the same, except that you also need to input \textcolor{RoyalBlue}{$Phi\_{}aprimeKron$} and \textcolor{RoyalBlue}{$Case2\_{}Type$}
   in exactly the same way as for the $ValueFnIter\_{}Case2$ command.
 \item By setting $n\_{}p$ to a non-zero value you tell the command that instead of using $fminsearch$ (which uses simplex methods) to solve
   for the general equilibrium using the general equilibrium conditions it should instead use a grid on prices. $n_p$ should describe the size of
   the grids for the price vectors (in exactly the same way as $n_d$ contains the size of the grids for the vector of decision variables $d$).
   The price grids themselves should be passed as a stacked column vector in $heteroagentoptions.pgrid$, the format of being exactly the same
   as is used for $d\_{}grid$ or $z\_{}grid$, ie. a stacked column vector. In this case the output of $p\_{}eqm\_{}index$ is now the grid
   index corresponding to $p\_{}eqm$, and $GeneralEqmConditions$ contains the values for each of the general equilibrium conditions at every point
   on the price grid.
\end{itemize}

\subsection{Options}
The options for the value function and stationary distribution can be set using structures called \textit{vfoptions}
and \textit{simoptions} and are exactly the same as those for $ValueFnIter\_{}Case1$ and $StationaryDist\_{}Case1$.
Further options can be passed in \textit{heteroagentoptions} and include $heteroagentoptions.pgrid$, \\
$heteroagentoptions.verbose$, $heteroagentoptions.multiGEcriterion$ and $heteroagentoptions.fminalgo$.
$heteroagentoptions.pgrid$ is used when you want to evaluate the general equilibrium conditions on a grid on prices, see
'Some Further Remarks' just above. $heteroagentoptions.multiGEcriterion$ and $heteroagentoptions.fminalgo$
allow you to change multiple general equilibrium conditions are evaluated (the former) and which algorithm is used
to compute the general equilibrium fixed-point problem (the later); both are currently inactive and simply implemented
to allow future extensions of the code.

\subsection{Examples}
See the Aiyagari (1994) example at \\
\href{https://github.com/vfitoolkit/VFItoolkit-matlab-examples/tree/master/HeterogeneousAgentModels}{github.com/vfitoolkit/VFItoolkit-matlab-examples/tree/master/HeterogeneousAgentModels}


